\section{Related Work}

\subsection{Active Range Sensing vs. Monocular Vision}
Planar 2D LiDARs have served as the benchmark sensor for mobile robot mapping and navigation stacks such as ROS 2 Nav2~\cite{macenski2020marathon, macenski2022robot}. Despite their reliability, mechanical LiDARs introduce mechanical wear, payload penalties, and significant cost bottlenecks. Solid-state LiDARs and RGB-D cameras partially address mechanical degradation, yet face stringent field-of-view (FOV) limits and ambient illumination failures~\cite{liu2024review}. Consequently, reconstructing range data from standard passive monocular cameras has emerged as a compelling alternative for lightweight autonomous platforms.

\subsection{Geometric Vision and Depth Estimation in Robotics}
Extracting metric spatial structures from monocular images requires geometric constraints or data-driven priors. In structured indoor environments, geometric properties such as vanishing points and ground planes enable precise parameter estimation. For instance, Lee and Hwang~\cite{lee2020fast} developed a robust self-calibration algorithm utilizing vanishing point detection in manmade indoor structures, demonstrating that structured architectural features provide stable geometric cues. Furthermore, combining sparse depth cues with neural representations was explored in~\cite{lee2020real}, where recurrent CNNs achieved real-time depth estimation for visual SLAM systems. In long-term robot operation, visual place recognition under dynamic environmental variations was investigated in~\cite{lee2019bag}, highlighting the importance of robust visual abstractions. Building upon these geometric and visual perception foundations, the proposed V-LiDAR leverages ground-plane homography to achieve deterministic $O(1)$ metric range projection without requiring end-to-end depth estimation inference.

\subsection{Floor Segmentation and Virtual Scan Generation}
Floor segmentation isolates traversable ground regions from surrounding obstacles. Traditional computer vision relied on color thresholds, edge gradients, and planar homographies. Recent deep convolutional networks and lightweight vision models, such as YOLOv11-seg~\cite{ultralytics2024yolo11}, offer robust pixel-wise semantic segmentation under complex textural variations. By treating the lower boundary of non-floor regions as the physical intersection between obstacles and the ground floor, virtual scan lines can be synthesized. However, previous works rarely integrated such virtual scans directly into production ROS 2 Nav2 costmaps with rigorous multi-run empirical validation and optical blind-spot analysis.
